\section{Detailed Design}
\label{sec:design}

We detail \sys{}'s unfenced credit signaling and ordering safety
(\S\ref{sec:design-signaling}), PCIe/mesh interference boundary
(\S\ref{sec:design-cliff}), sizing bounds (\S\ref{sec:design-sizing}), and
inline multi-core buffering (\S\ref{sec:design-fifo}), summarized in
Algorithm~\ref{alg:creditring}.

\subsection{Unfenced 64-Byte MMIO Credit Signaling}
\label{sec:design-signaling}

Each host queue $q$ owns an exclusive, 64-byte-aligned BAR4 cache line
\texttt{CreditLine[q]} (512\,B total for 8 queues, zero locks or false sharing)
containing \texttt{[gen\_head, taken, run\_id, ..., gen\_tail]}. Once queue $q$
consumes and refills at least $B$ packets, the host publishes \texttt{taken}
with a single \textbf{unfenced} \texttt{MOVDIR64B} posted write---no
\texttt{SFENCE}.

\paragraph{Why \texttt{SFENCE} is unnecessary.}
Unlike multi-line payload writes that require fencing payloads before ringing a
doorbell~\cite{liew2026asplos}, a credit update is a single 64-byte atomic
monotonic counter:
\begin{enumerate}
\item \emph{Atomic same-address ordering:} \texttt{MOVDIR64B} emits one atomic
  64-byte posted PCIe write, and PCIe preserves ordering among posted writes to
  the same address.
\item \emph{Monotonic idempotence with ARMv8 seqlock:} The SmartNIC reads
  \texttt{CreditLine[q]} from local coherent L3 using acquire-ordered loads and
  checks \texttt{gen\_head == gen\_tail}, retrying on concurrent invalidation
  ($\sim$28.7\% under tight polling, 0 torn reads across 618M+ packets), then
  applies \texttt{taken\_cached = max(taken\_cached, line.taken)}. Coalesced
  or overwritten intermediate updates are therefore harmless.
\item \emph{Ordering safety for static and dynamic host rings:} On static
  in-place rings, descriptors never change after startup, so no write-back (WB)
  descriptor ordering dependency exists. On standard dynamic \texttt{mbuf}
  replacement, the DPDK driver writes new descriptors and rings the uncached
  hardware doorbell (\texttt{SLIST\_DBELL}) every $B_{\text{refill}}$ packets
  (64 on \texttt{octeon\_ep}). Even if an unfenced \texttt{MOVDIR64B} arrives
  slightly ahead of a concurrent refill doorbell, hardware descriptor readiness
  remains enforced by \texttt{SLIST\_DBELL}, while premature watermark drops
  are precluded whenever
  \begin{equation}
    R_{\text{eff}} + B_{\text{refill}} \;\le\; D - W_{\text{mark}}.
    \label{eq:wmark-invariant}
  \end{equation}
  For $D=512$, $W_{\text{mark}}=128$, and $B_{\text{refill}}=64$, any
  $R_{\text{eff}} \le 320$ (we use 128--256) stays strictly safe even when
  credit and refill boundaries are unaligned.
\end{enumerate}

\paragraph{Single-write-per-burst rule.}
If the host finishes a 64-packet receive burst and loops $64/B$ times emitting
back-to-back writes to the \emph{same} BAR4 cache line (e.g., 8 writes at
$B=8$), 85\% of intermediate writes are overwritten before the SmartNIC polls
while triggering same-line mesh serialization (\S\ref{sec:eval-e6}). \sys{}
therefore issues \emph{at most one} \texttt{MOVDIR64B} carrying the latest
\texttt{taken} per receive burst whenever $\text{taken} - \text{published} \ge B$,
backed by an idle-poll timeout ($\sim$50\,$\mu$s or 64 empty polls) so tail
packets never strand credits.

\begin{algorithm}[t]
\footnotesize
\noindent\rule{\linewidth}{0.6pt}\\[2pt]
\textbf{SmartNIC Core (per queue $q$, inline on-path loop):}
\begin{enumerate}[leftmargin=1.4em, itemsep=0.5pt, topsep=2pt]
  \item \texttt{n = eth\_rx\_burst(mac\_port, q, pkts, 64)}
  \item \textbf{if} \texttt{n > 0}: append \texttt{pkts[0..n)} to \texttt{elastic\_fifo[q]}
  \item \texttt{taken\_q = max(taken\_q, seqlock\_read(BAR4\_line[q]))}
  \item \texttt{avail = R\_eff - (submitted\_q - taken\_q)}
  \item \textbf{if} \texttt{avail > 0} \textbf{and} \texttt{!elastic\_fifo[q].empty()}:
  \item \quad \texttt{k = min(avail, elastic\_fifo[q].len(), 64)}; pop \texttt{k} pkts
  \item \quad \texttt{sdp\_tx\_burst(q, k)}; \texttt{submitted\_q += k}
\end{enumerate}
\vspace{3pt}
\textbf{Host Core (per queue $q$, unmodified DPDK RX + 1-line credit hook):}
\begin{enumerate}[leftmargin=1.4em, itemsep=0.5pt, topsep=2pt]
  \item \texttt{n = rte\_eth\_rx\_burst(sdp\_port, q, mbufs, 64)} \hfill\textit{// normal DMA}
  \item process \texttt{mbufs[0..n)} in unmodified app; \texttt{rte\_pktmbuf\_free\_bulk}
  \item \texttt{taken\_q += n}
  \item \textbf{if} \texttt{taken\_q - published\_q >= B} \textbf{or} \texttt{idle\_flush()}:
  \item \quad \texttt{movdir64b\_nofence(BAR4\_line[q], \{++gen, taken\_q, ..., gen\})}
  \item \quad \texttt{published\_q = taken\_q} \hfill\textit{// at most 1 write per burst}
\end{enumerate}
\vspace{-2pt}
\noindent\rule{\linewidth}{0.6pt}
\caption{Unified single-path credit gating on the SmartNIC and unfenced credit
return on the host.}
\label{alg:creditring}
\end{algorithm}

\subsection{The Unfenced MMIO Interference Cliff}
\label{sec:design-cliff}

High-rate unfenced MMIO is unsafe for bulk data transfers~\cite{vuppalapati2024hostnet}:
as we prove in \S\ref{sec:eval-e1}, driving one Sapphire Rapids core at
unconstrained $\sim$175\,Mpps saturates the PCIe inbound write-completion
queue past a sharp \textbf{$>$94\% busy cliff} (98.9\%), inflating co-located
L3 latency by +168.5\% and cutting 6-core STREAM bandwidth by 25.0\% even as
aggregate mesh occupancy drops (387.6 $\to$ 268.5) due to injection
backpressure. Restricting unfenced MMIO strictly to batched credits keeps
aggregate traffic within \textbf{0.5--7.0\,Mpps} across 1--8 cores (PCIe ingress
$\le 0.3\%$, mesh occupancy $\le 2.8$), 25--350$\times$ below the cliff with
$\le 0.35\%$ co-located neighbor impact.
\subsection{Sizing Rings, Pools, and Batch Granularity}
\label{sec:design-sizing}

Because standard DMA binds one descriptor and one \texttt{mbuf} per packet,
the minimum lossless credit window depends on \emph{packet rate} rather than
raw bytes:
\begin{equation}
  R_{\min} = \lceil \lambda_{\text{pps}} \cdot \tau_{\text{rtt}} \rceil + B.
  \label{eq:bdp-slots}
\end{equation}
With $\tau_{\text{rtt}} = 2.15$--$2.53\,\mu\text{s}$ (\texttt{ism\_enable=1},
$B=32$), 1500\,B at 100GbE (8.22\,Mpps) needs only $\sim$53 slots (77.7\,KiB
wire bytes; 106\,KiB of 2\,KiB \texttt{mbuf}s), whereas 64\,B at 90--148.8\,Mpps
requires 226--353 descriptor slots (452--706\,KiB of 2\,KiB \texttt{mbuf}s).
Even under conservative non-ISM polling ($\tau_{\text{rtt}} \approx 7.6$--$8.4\,\mu\text{s}$),
$R_{\text{eff}} = 128$--$256$ clears the 96-slot 1500\,B knee comfortably.

To prevent cold-memory sweeping, tight credits must be paired with a compact
host \texttt{mempool} ($\le 10$\,MiB aggregate, $\le 7.5$\,MiB actual 1500\,B
DMA write footprint), and return batching must satisfy
\begin{equation}
  B \;\le\; R_{\text{eff}} / 2
  \label{eq:batch-bound}
\end{equation}
so the SmartNIC overlaps DMA submission with host credit return without
starving ($B=32$--$64$ for $R_{\text{eff}}=128$--$256$, $<0.2\%$ reverse PCIe
bandwidth at 1500\,B).

\subsection{Elastic Buffering and Multi-Core Isolation}
\label{sec:design-fifo}

\sys{} exploits the SmartNIC's two-tier memory (48\,MiB L3 + 6-channel DDR5)
using either a head-resetting \texttt{CompactSlab} FIFO---resetting offset to 0
on empty so sub-L3 bursts remain hot and small frames avoid 9$\times$
\texttt{mbuf} space amplification---or a zero-copy \texttt{MbufRing} backed by
hardware \texttt{NPA} LIFO pools to avoid ARM copies on 1500\,B frames. Under
sustained single-queue overload, capping \texttt{MbufRing} depth or selecting
\texttt{CompactSlab} prevents perpetually full queues from thrashing SmartNIC
L3 (\S\ref{sec:eval-e4}). Across $N$ host cores, each RSS queue maps 1-to-1 to
an isolated SmartNIC FIFO, \texttt{NPA} pool, DMA queue, 64\,B BAR4 line, and
compact host pool, isolating 1--10\,ms core stalls without locks. Finally, the
SmartNIC gate requires only 16\,B of state and an $O(1)$ counter comparison per
queue, making it directly amenable to fixed-function NIC DMA hardware offload.
